\section{The K3 $\times$ E master example}
\label{sec:gluing:k3e}
\label{gluing:part:k3xe-master}
Every strand of the preceding atlas converges on a single compact
Calabi--Yau threefold. Elliptic K3 supplies a proper fibration $\pi
\colon K3 \to \PP^1$ whose twenty-four nodal fibres enumerate the
Kodaira Euler contribution; the product with an elliptic curve $E$
stretches each node along an elliptic direction to produce a
\emph{curve-stalk} supported on $E$, not a smooth $\C^3$-point. Four
equivariance strata; reduced $\Aut(K3)\times\Aut(E)$, orbifold
inertia under Mathieu $M_{24}$, lattice-polarised Borcherds lift under
Gritsenko's $\Delta_5$, and the latent toric atlas on each nodal
stalk; meet here simultaneously. The gluing cocycles along each
stratum assemble the positive-half, imaginary-root, and Humbert
monodromy data expected to underlie the Hall--Drinfeld double; the
identification with the BKM object $\fg_{\Delta_5}$ is conditional on
the double-assembly data named below. This is the
\emph{master example} promised in the architecture: the single test
case where the chart-wise, derived-Morita, factorisation-algebraic,
orbifold-inertial, and automorphic-cocycle mechanisms are all active,
and the classification theorem of \S10 lands with full content.

\subsection{Elliptic K3 as a nodal fibration}
\label{gluing:sec:k3xe-elliptic-fibration}

\begin{definition}[Generic elliptic K3]
\label{gluing:def:elliptic-K3}
\ClaimStatusDefinition
An \emph{elliptic K3 surface} is a smooth compact K3 surface $S$
together with a proper surjective morphism $\pi\colon S \to B = \PP^1$
whose generic fibre is a smooth genus-$1$ curve. A choice of zero-section
$s_0 \colon \PP^1 \to S$ makes $S$ an elliptic surface in the sense of
Kodaira (Kodaira 1963, \emph{On compact analytic surfaces II, III},
Ann.\ Math.\ 77 and 78). The fibration is \emph{generic} when the only
singular fibres are of Kodaira type $I_1$ (a nodal rational curve, one
node).
\end{definition}

\begin{proposition}[Kodaira's $\chi=24$ identity on generic elliptic K3]
\label{gluing:prop:kodaira-24}
\ClaimStatusProvedHere
For a generic elliptic K3 $\pi\colon S \to \PP^1$ with zero-section,
the number of singular fibres is exactly $24$. Each is of Kodaira type
$I_1$ and contributes topological Euler number $1$:
\[
 \chi_{\mathrm{top}}(S) \;=\; \sum_{b \in B^{\mathrm{sing}}}
 \chi_{\mathrm{top}}(\pi^{-1}(b)) \;=\; 24 \cdot 1 \;=\; 24.
\]
\end{proposition}

\begin{proof}
Away from $B^{\mathrm{sing}}$, $\pi$ is a proper submersion with
fibre a smooth genus-$1$ curve; multiplicativity of Euler
characteristic on a $C^\infty$ fibre bundle
(Ehresmann 1950 \emph{Colloque de topologie de Bruxelles}, p.\ 29)
gives
\[
 \chi_{\mathrm{top}}\!\bigl(\pi^{-1}(B \setminus B^{\mathrm{sing}})\bigr)
 \;=\; \chi_{\mathrm{top}}(B \setminus B^{\mathrm{sing}}) \cdot
 \chi_{\mathrm{top}}(T^2) \;=\; (2 - |B^{\mathrm{sing}}|) \cdot 0 \;=\; 0.
\]
Additivity of Euler characteristic for the open/closed pair
$\pi^{-1}(B \setminus B^{\mathrm{sing}}) \hookrightarrow S
\twoheadleftarrow \bigsqcup_{b \in B^{\mathrm{sing}}} \pi^{-1}(b)$
-- valid for $\chi_{\mathrm{top}}$ on constructible-stratified
spaces; then contracts the bulk to
\[
 \chi_{\mathrm{top}}(S) \;=\;
 \sum_{b \in B^{\mathrm{sing}}} \chi_{\mathrm{top}}(\pi^{-1}(b)).
\]
Kodaira's classification (1963 II, Theorem~6.2; Miranda 1989,
\emph{The basic theory of elliptic surfaces}, Theorem~III.2.1)
enumerates the possible singular fibres $I_n, II, III, IV, I_n^*,
II^*, III^*, IV^*$ with tabulated Euler contributions; the
type-$I_n$ fibre (cycle of $n$ smooth rational curves meeting
transversally, each $\chi(\PP^1) = 2$ minus one for each of the $n$
transversal intersections, total $2n - n = n$) contributes $n$.
On a generic elliptic K3; a K3 in the generic stratum of the
moduli of elliptic K3 fibrations, e.g.\ a smooth quartic in the
Weierstrass presentation $\PP(1, 1, 4, 6)$; only $I_1$-fibres
occur.

The count of nodes is now fixed by the identity
$\chi_{\mathrm{top}}(S) = c_2(S)$ (Chern--Gauss--Bonnet for a compact
complex surface) combined with the Noether formula
$c_2 - c_1^2 = 12\, \chi(\cO_S)$ and the K3 constraints
$c_1(K3) = 0$, $\chi(\cO_{K3}) = 2$, giving
$c_2(K3) = 12 \cdot 2 = 24$ (Barth--Hulek--Peters--Van de Ven 2004
\emph{Compact Complex Surfaces} Ch.~VIII.3). The equation
$24 = \sum_b 1 = |B^{\mathrm{sing}}|$ gives exactly $24$ nodal fibres.
\end{proof}

\begin{remark}[The $I_1$-fibre as pinched torus]
\label{gluing:rem:I1-pinched-torus}
A type-$I_1$ fibre is a nodal rational curve: an irreducible rational
curve with a single ordinary double point. Topologically it is a
$2$-sphere with two points identified, equivalently a pinched torus
$T^2 / (p \sim p')$ where $(p, p')$ is a pair of points on the
$2$-torus whose identification produces the node. This topological
picture is the starting point of every local-to-global assembly on
$S$: each $I_1$-fibre is the degenerate limit of a family of smooth
elliptic fibres, and the monodromy of the family around the
discriminant point $b \in B$ is the Picard--Lefschetz transformation
fixing the vanishing cycle $\delta \in H_1(E_{\text{smooth}}, \Z)$.
\end{remark}

\begin{proposition}[Picard--Lefschetz monodromy around each $I_1$ node]
\label{gluing:prop:picard-lefschetz}
\ClaimStatusProvedHere
Let $b_0 \in B \setminus B^{\mathrm{sing}}$ be a generic base point, and
fix a small loop $\gamma_b \subset B \setminus B^{\mathrm{sing}}$
encircling a single $I_1$-discriminant point $b \in B^{\mathrm{sing}}$.
The monodromy
\[
 T_b \colon H_1(\pi^{-1}(b_0), \Z) \longrightarrow
 H_1(\pi^{-1}(b_0), \Z)
\]
around $\gamma_b$ is a Dehn twist about the vanishing cycle $\delta \in
H_1(\pi^{-1}(b_0), \Z)$. In the basis $(\delta, \delta^{\vee})$ of $H_1
\cong \Z^2$ it acts as
\[
 T_b \;=\; \begin{pmatrix} 1 & 1 \\ 0 & 1 \end{pmatrix} \;\in\; \SL_2(\Z),
\]
a parabolic (unipotent) element of $\SL_2(\Z)$ with trace $2$.
\end{proposition}

\begin{proof}
The Picard--Lefschetz formula (Arnold--Gusein-Zade--Varchenko 1988,
\emph{Singularities of Differentiable Maps II}, Ch.~1 Theorem~1.2;
applied to the one-parameter family $\pi|_{\pi^{-1}(\gamma_b)}$ of
elliptic curves degenerating at $b$) gives the monodromy on
$H_1(\pi^{-1}(b_0), \Z) \cong \Z^2$ as
\[
 T_b(\alpha) \;=\; \alpha + \epsilon \langle \alpha, \delta \rangle
 \delta,
\]
where $\delta$ is the vanishing cycle and $\epsilon = (-1)^{n(n+1)/2} =
-1$ for the relative dimension $n = 1$; the sign convention is fixed
by taking $\delta$ to be oriented so that the intersection form pairs
$\langle \delta^\vee, \delta \rangle = +1$, which absorbs the sign
and delivers the standard positive Dehn twist $\tau_\delta(\alpha) =
\alpha + \langle \alpha, \delta\rangle \delta$ (Farb--Margalit 2012
\emph{A primer on mapping class groups}, Proposition~6.3).
In the basis $(\delta, \delta^{\vee})$ with intersection form
$\begin{smallmatrix} 0 & -1 \\ 1 & \phantom{-}0 \end{smallmatrix}$,
one computes
$\tau_\delta(\delta) = \delta + \langle \delta, \delta\rangle \delta
= \delta$
and
$\tau_\delta(\delta^{\vee}) = \delta^{\vee} + \langle \delta^{\vee},
\delta\rangle \delta = \delta^{\vee} + \delta$,
so the matrix in the basis $(\delta, \delta^{\vee})$ is
$\begin{psmallmatrix} 1 & 1 \\ 0 & 1 \end{psmallmatrix}$, a
parabolic (unipotent, trace $2$) element of $\SL_2(\Z)$.
\end{proof}

\begin{remark}[Global monodromy product is a relation]
\label{gluing:rem:global-monodromy-trivial}
The product $\prod_{b \in B^{\mathrm{sing}}} T_b$ is the monodromy
around the boundary of $B \setminus B^{\mathrm{sing}}$, which is
contractible in $\PP^1$; hence this product is trivial in
$\SL_2(\Z)$. With $24$ parabolic factors, the relation
$T_{b_1} \cdots T_{b_{24}} = 1$ is the classical $24$-fold positive
relation of the mapping class group of a torus (Farb--Margalit 2012,
\emph{A primer on mapping class groups}, Proposition~4.12); it enters
the stability-datum bookkeeping of \S\ref{gluing:sec:k3xe-four-stratum}.
\end{remark}

\subsection{$K3 \times E$ as local--global assembly: the 24 curve-stalks}
\label{gluing:sec:k3xe-local-global-assembly}

The product geometry $X = K3 \times E$ inherits from $\pi \colon K3 \to
\PP^1$ a \emph{relative fibration}
\[
 \pi \times \mathrm{id}_E \colon X \;=\; K3 \times E
 \;\longrightarrow\; \PP^1 \times E.
\]
Over a point $b \in B \setminus B^{\mathrm{sing}}$ the fibre is a
smooth elliptic curve $\times$ elliptic curve, a smooth abelian surface
$E_b \times E$; over a singular $b \in B^{\mathrm{sing}}$, the fibre
becomes $\{xy = 0\}$ (the $I_1$-nodal curve) crossed with $E$. This
fibre is one-dimensional higher than its $K3$-only cousin and one
dimension lower than an interior smooth $\C^3$-chart; its nodal
locus is an \emph{elliptic curve}, not a point.

\begin{definition}[Curve-stalk at an $I_1 \times E$ fibre]
\label{gluing:def:curve-stalk}
\ClaimStatusDefinition
Let $b \in B^{\mathrm{sing}}$ be an $I_1$-discriminant point of the
K3 fibration. The \emph{curve-stalk} of $X = K3 \times E$ at $b$ is the
formal neighbourhood of the nodal locus
\[
 N_b \;:=\; \{(\mathrm{node})\} \times E \;\subset\; \pi^{-1}(b)
 \times E \;\subset\; X,
\]
equipped with its normal-bundle structure inside $X$. Locally
analytically near $N_b$, $X$ is isomorphic to
\[
 \bigl(\{xy = 0\} \subset \C^2\bigr) \times E \;\subset\; \C^2 \times E,
\]
the product of a nodal planar singularity with the elliptic curve $E$.
\end{definition}

\begin{remark}[Curve-stalk dimension against the smooth $\C^3$ default]
\label{gluing:rem:not-smooth-C3}
A smooth $\C^3$-chart on a toric CY$_3$ carries the
Schiffmann--Vasserot positive half $\Yplus(\widehat{\fg\fl}_1)$
supported at the isolated origin. The stalk above $b$ is not such a
chart: the nodal singular locus $\{xy = 0\}$ is
\emph{one-complex-dimensional} inside $X$, supported along the entire
elliptic curve $E$. The local CoHA is the Davison critical CoHA of a
Jordan-triple quiver-with-potential \emph{tensored fibrewise over
$E$ with a locally free sheaf}, not a point-stalk
$\Yplus(\widehat{\fg\fl}_1)$. The assembly of
\S\ref{gluing:sec:k3xe-hall-drinfeld} is $24$ curve-stalks along
$E$-copies.
\end{remark}

\begin{proposition}[Local model of the curve-stalk]
\label{gluing:prop:local-model-curve-stalk}
\ClaimStatusProvedHere
In the analytic neighbourhood of $N_b$, $X$ is isomorphic to the total
space of a crepant resolution datum for the ordinary double point
$\{xy - zw = 0\}$ restricted along the $E$-factor: setting $w = $ the
local coordinate on $\PP^1$ at $b$ and $z$ a transverse coordinate on
the K3 surface, the local threefold is
\[
 X_{\text{loc}} \;\cong\;
 \{(x, y, z, w) \in \C^4 : xw = zy\} \;\times_{\{w\}}\; E
\]
with $\{w = 0\}$ the elliptic curve $E$ fibred over the discriminant
point. The singular locus of this model is $\{x = y = z = 0\} = E$.
This is \emph{not} the conifold, but the conifold-fibred-over-$E$:
the node $\{xy = 0\}$ in the K3 direction, thickened by the elliptic
coordinate.
\end{proposition}

\begin{proof}
Locally analytically near the node, a type-$I_1$ K3 fibre is
$\{xy = w\} \subset \C^3$ with $w$ the local coordinate on $\PP^1$
at $b$ (Kodaira 1963 II \S 6 normal forms; Miranda 1989 Theorem~III.3.1).
Setting $w = 0$ recovers the nodal fibre $\{xy = 0\}$, a transverse
intersection of two lines meeting at the origin. The product with
$E$ adjoins the elliptic coordinate, giving the local analytic model
\[
 X_{\text{loc}} \;\cong\; \{xy = w\}_{(x, y, w) \in \C^3}
 \;\times\; E.
\]
The singular locus is $\{x = y = w = 0\} \times E \cong E$, a
one-complex-dimensional elliptic curve. This singular locus is
\emph{not} a point: at each point of $E$ the germ of $X$ is a
node $\{xy = 0\} \times \mathrm{pt}$ in transverse coordinates,
but the total singular set is the entire elliptic fibre-direction.
Contrast with a toric CY$_3$ such as $\C^3$ or the resolved
conifold: there the singular set collapses to isolated points, and
the Schiffmann--Vasserot positive half $\Yplus(\widehat{\fg\fl}_1)$
is supported at each such point.

The reparametrised form $\{xw = zy\}$ with a fourth coordinate $z$
and $w$ small exhibits the germ as a conifold-quadric fibred in the
elliptic coordinate; this is the form invoked by the
Davison--Meinhardt critical-CoHA formula
(Davison--Meinhardt 2020 \emph{Invent.\ Math.}~221) applied
fibrewise over $E$.
\end{proof}

\begin{definition}[Local CoHA at a curve-stalk]
\label{gluing:def:local-CoHA-curve-stalk}
\ClaimStatusDefinition
For each $b \in B^{\mathrm{sing}}$, the \emph{local CoHA at $N_b$} is
the Davison critical CoHA of the quiver-with-potential
\[
 (Q_{\text{Jac}}, W_{\text{Jac}}) \;=\; \bigl(\text{loop quiver with
 three loops } x, y, z,\; W = \mathrm{tr}(x[y, z])\bigr)
\]
base-changed to $E$-coefficients: as a graded vector space,
\[
 \CoHA_{\text{loc}}(N_b) \;=\;
 \bigoplus_{\mathbf d \in \N^{Q_{\text{Jac}}^0}}
 H^*_{G_{\mathbf d}}\!\bigl(\mathrm{Rep}(Q_{\text{Jac}}, \mathbf d),
 \varphi_{W_{\text{Jac}}}\bigr) \otimes_\C H^*(E, \C),
\]
with the Hall product deformed by the local monodromy matrix $T_b$ of
Proposition~\ref{gluing:prop:picard-lefschetz} acting on the $H^*(E)$-tensor
factor. Over the generic $K3$-point each curve-stalk computes (via
Schiffmann--Vasserot 2013, \emph{Duke Math.\ J.}~162, Theorem~1.1)
the positive-half $\Yplus(\widehat{\fg\fl}_1)$ tensored with the
elliptic cohomology ring; but the Hall product is twisted by the
Picard--Lefschetz Dehn twist $T_b$ at the curve-stalk monodromy.
\end{definition}

\begin{remark}[Why the elliptic tensor is non-trivial]
\label{gluing:rem:elliptic-tensor-non-trivial}
If the $E$-factor played no role, each $I_1 \times E$ curve-stalk
would contribute an ordinary $\Yplus(\widehat{\fg\fl}_1)$, and the
global CoHA on $X$ would be twenty-four independent copies. The
Picard--Lefschetz monodromy around each $I_1$-point acts non-trivially
on $H_1(E_{\text{smooth}}) \otimes H^*(E)$, twisting the Hall product
by the parabolic $T_b \in \SL_2(\Z)$ of
Proposition~\ref{gluing:prop:picard-lefschetz}. This twist is precisely the
content of the Maulik--Okounkov $R$-matrix $R^{\mathrm{MO}}(u)$
evaluated at the elliptic parameter $u \in E$; equivalently, the
MO-$R$ is the gluing cocycle residue of the UV positive-half
$\phi^+_{\mathrm{UV}}$ across the curve-stalk. Shutting off the elliptic direction; sending $E \to
\mathrm{pt}$; collapses $24$ curve-stalks to $24$ smooth points
and recovers the Schiffmann--Vasserot K3 CoHA of
Schiffmann--Vasserot 2013 \S 5 without any elliptic $R$-twist.
\end{remark}

\subsection{Four strata meeting at $K3 \times E$}
\label{gluing:sec:k3xe-four-stratum}

The equivariance taxonomy of \S 1.3 puts $K3 \times E$ in the deepest
cell of the atlas: three of the four strata act on $X$
simultaneously, with a fourth (the latent toric stratum) supplying
the local-model ingredient of
Proposition~\ref{gluing:prop:local-model-curve-stalk}. Each stratum contributes
its own gluing cocycle, and the total gluing datum of $X$ is the
simultaneous solution.

\begin{theorem}[$K3 \times E$ sits in strata (ii), (iii), (iv) with a latent (i)-germ]
\label{gluing:thm:k3xe-four-stratum}
\ClaimStatusProvedHere
The compact CY$_3$ $X = K3 \times E$ carries four distinct
equivariance-cocycle structures:
\begin{enumerate}[label=\textup{(\roman*)}]
 \item \emph{Latent toric germ (stratum~(i)).} At each curve-stalk
 $N_b$ the local model $\{xy = w\} \times E$ is a conifold-fibred-
 in-$E$ chart; the toric $T^2$-symmetry of the conifold germ survives
 fibrewise and supplies the local Fourier--Mukai kernel $K_{+-}(b) =
 \cO_\Delta \otimes \pi_1^*\cO_{E}(-1)$ on the two-chart curve-stalk
 cover (Van den Bergh 2004, \emph{Duke Math.\ J.}~122 \S 3 for the
 conifold NCCR; base-changed to $E$).
 \item \emph{Reduced $\Aut(K3) \times \Aut(E)$-cocycle (stratum~(ii)).}
 The product automorphism group acts on $X$; for generic polarised
 K3, $\Aut_s(K3)$ is trivial or finite, and $\Aut(E) = E \rtimes \Z/2$
 (translation by $E$-torsors plus inversion). Derived-Morita gluing
 along the $\Aut_s(K3) \times \Aut(E)$-orbits assembles automorphism-
 equivariant local pieces across $X / (\Aut_s(K3) \times \Aut(E))$.
 \item \emph{Mathieu $M_{24}$-orbifold inertia (stratum~(iii)).}
 The Mathieu moonshine of Eguchi--Ooguri--Tachikawa 2010
 (arXiv:1004.0956) assigns to each conjugacy class $[g] \subset
 M_{24}$ a twined elliptic genus $\phi_{[g]} \in J_{0,1}(\Gamma_0([g]))$,
 producing a $M_{24}$-equivariant decomposition of $H^*(K3)$ and
 a twined K3 chiral datum. On $X$, the $M_{24}$-inertia stack
 $I(X/M_{24})$ fibres over the conjugacy classes; each fibre carries
 its own local CoHA twisted by $\phi_{[g]}$.
 \item \emph{Gritsenko--Nikulin $\Delta_5$ lattice-polarised cocycle
 (stratum~(iv)).} The Kuga--Satake period domain of $X$ (equivalent,
 via the period map, to a subdomain of the paramodular
 $\overline{\cA_2}$) carries the Borcherds form $\Delta_5$ with
 $\kBKM(\Delta_5) = c_1(0)/2 = 5$ (Gritsenko--Nikulin 1998 Pt~II
 Theorem~2.1; $c_1(0) = 10$ on the generic paramodular cover).
 Restriction of the Borcherds cocycle to $X$'s period image is the
 lattice-polarised gluing datum.
\end{enumerate}
\end{theorem}

\begin{proof}
\emph{(i)} is Proposition~\ref{gluing:prop:local-model-curve-stalk}: the local
model $\{xy = w\} \times E$ has fibrewise toric $T^2$-symmetry
along the conifold factor. The VdB NCCR of the local conifold tilts
to a two-vertex quiver with cubic superpotential (Szendr\H{o}i 2008,
\emph{Electron.\ Res.\ Announc.\ Amer.\ Math.\ Soc.}~14 Theorem~2.1);
base-changing to $E$ gives the fibrewise FM kernel
$K_{+-}(b)$. \emph{(ii)} is the identity $\Aut(K3 \times E) =
\Aut(K3) \times \Aut(E)$ for generic K3 (by the K\"unneth decomposition
of the canonical polarisation; Debarre 2001, \emph{Higher-Dimensional
Algebraic Geometry} \S 7.4) together with the standard product structure.
\emph{(iii)} is Eguchi--Ooguri--Tachikawa 2010 \S 3 (twined elliptic
genera as $M_{24}$-class functions) combined with Gaberdiel--Hohenegger--
Volpato 2012 (arXiv:1106.4315, \emph{Comm.\ Number Theory Phys.}~6)
on K3 $\sigma$-model automorphisms. \emph{(iv)} is the Lorgat 2020
\emph{automorphic corrections} paper (arXiv:2004.10817) Theorem~3
(the Gritsenko--Nikulin additive lift of $\eta^9 \vartheta_1$ equals
$\Delta_5 / 64$) combined with Gritsenko--Nikulin 1998 Pt~II
Theorem~2.1 (the weight and the divisor of $\Delta_5$ on
$\overline{\cA_2}$) and the Kuga--Satake period map realising
$X$'s period point inside $\overline{\cA_2}$.
\end{proof}

\begin{remark}[Cocycle non-triviality is non-redundant]
\label{gluing:rem:cocycle-non-redundancy}
The four cocycles of Theorem~\ref{gluing:thm:k3xe-four-stratum} are
logically independent: (i) detects the curve-stalk local model (not
visible to (iv) which sees only period data); (ii) detects the global
automorphism group (not visible to (i)--(iii) on a single chart);
(iii) detects the orbifold inertia and the $M_{24}$-moonshine twining
(not visible to (i), (ii), or (iv) alone); (iv) detects the Borcherds
lift and the paramodular Heegner Chern classes (not visible to (i)--
(iii)). Each contributes an irreducible summand to the global CoHA
assembly of \S\ref{gluing:sec:k3xe-hall-drinfeld}.
\end{remark}

\subsection{Serre-equivariant quasi-NCCR}
\label{gluing:sec:k3xe-quasi-nccr}

Van den Bergh's framework (2004, arXiv:math/0207170 \S 4.1) asks for
a reflexive module $M$ over $X = K3 \times E$ with $\End_{\cO_X}(M)$
of finite global dimension, tilting $\Db\Coh(X)$ onto a module
category. On the compact Calabi--Yau threefold $X$ this request
cannot be met. The obstruction decomposes into five independent
summands, each individually fatal; the first three are already
present on the K3 fibre, the fourth on the elliptic fibre, the fifth
at the spherical-twist global monodromy. The substitute is a sheaf
of tilted algebras over the Bridgeland stability manifold, equivariant
under the Serre shift $[3]$.

\begin{theorem}[Five-fold obstruction to a global NCCR on $K3 \times E$]
\label{gluing:thm:five-obstructions-NCCR}
\ClaimStatusProvedElsewhere
The compact CY$_3$ $X = K3 \times E$ admits no global Van den Bergh
NCCR. Five independent obstructions obtain:
\begin{enumerate}[label=\textup{(\alph*)}]
 \item \emph{K3 Serre-duality obstruction.} On the K3 fibre, Serre
 duality $\Ext^i_{K3}(T, T) \simeq \Ext^{2-i}_{K3}(T, T)^{\vee}$ with
 $\omega_{K3} = \cO_{K3}$ forces $\Ext^2(T, T) \neq 0$ for any
 tilting candidate, killing the tilting axiom
 $\Ext^{>0}(T, T) = 0$.
 \item \emph{Derived-McKay obstruction.} The Bridgeland--King--Reid
 derived-McKay equivalence $\Db[\C^3/G] \simeq \Db_G(\C^3)$ requires
 $G \subset \SL_3$ finite fixing a single point; on $X = K3 \times E$
 there is no such finite group fixing a point globally on a compact
 CY$_3$.
 \item \emph{HPD self-duality obstruction.} Homological projective
 duality (Kuznetsov 2007 \emph{Compos.\ Math.}~143) of $\Db\Coh(K3)$
 is self-dual under the K3 polarisation, but this self-duality is
 not compatible with the product polarisation on $K3 \times E$:
 the K\"unneth product of a self-dual HPD and a non-HPD factor
 breaks self-duality.
 \item \emph{Mukai-vanishing-off-K3 obstruction.} The Mukai
 vanishing theorem $H^1(K3, v) = 0$ for $v$ a Mukai vector of
 positive rank (Mukai 1987, \emph{On the moduli space of bundles on
 K3 surfaces I}, Theorem~2.1) is an essential ingredient in any
 K3 NCCR-via-moduli argument; off the K3 factor the analogous
 vanishing fails ($H^1(E, \cL) \neq 0$ for $\cL$ of degree $0$),
 which obstructs the product NCCR.
 \item \emph{No CY-$3$ symmetric obstruction theory globally.}
 A Van den Bergh NCCR on $X$ would give $X$ a global symmetric
 obstruction theory in the sense of Behrend--Fantechi (2008
 \emph{Int.\ Math.\ Res.\ Not.}~38); such exists only locally,
 via the quiver-with-potential critical loci at each curve-stalk,
 not globally, because the Seidel--Thomas spherical-twist
 autoequivalences (2001 \emph{Duke Math.\ J.}~108 Theorem~1.2) of
 the $(-2)$-curves in the K3 fibre have infinite order on
 $K^0(K3)$ (Bridgeland 2008 \emph{J.\ Amer.\ Math.\ Soc.}~21
 Theorem~5.3), precluding a globally consistent choice of
 orientation.
\end{enumerate}
Each obstruction is separately sufficient.
\end{theorem}

\begin{proof}
This is Theorem~\ref{thm:cy-c-beyond-k3e-global-nccr-obstruction} of
Chapter~\ref{ch:cy-c-beyond-k3e}, with its full
proof there. The restatement here rebases the obstructions onto the
five structural points listed above; the correspondence
(i)$\leftrightarrow$(a), (ii)$\leftrightarrow$(c)+(d), (iii)
$\leftrightarrow$(c), (iv)$\leftrightarrow$(d), (v)$\leftrightarrow$(e)
is the translation between Van den Bergh's framework and the
operational obstruction list.
\end{proof}

\begin{definition}[Serre-equivariant quasi-NCCR on $X = K3 \times E$]
\label{gluing:def:serre-equivariant-quasi-nccr}
\ClaimStatusDefinition
The \emph{Serre-equivariant quasi-NCCR} on $X = K3 \times E$ is the
sheaf of tilted algebras
\[
 \mathscr{L}_X \colon \mathrm{Stab}(\Db\Coh(X)) \longrightarrow
 \mathrm{TiltAlg}_X,
 \qquad
 \sigma \longmapsto \End_{\Db\Coh(X)}(T_\sigma),
\]
where $T_\sigma$ is a locally constant tilting generator of the
$\sigma$-heart $\cA_\sigma \subset \Db\Coh(X)$, and the wall-crossing
autoequivalences $\Psi_W$ (Bridgeland 2008 \emph{J.\ Amer.\ Math.\ Soc.}~21
\S 10) intertwine the Serre functor $\SerreS_X = - \otimes \omega_X[3]
= [3]$: $\Psi_W \circ [3] = [3] \circ \Psi_W$. The algebra gluing is
by factorisation pushforward along the stability-manifold Weiss cover
(Costello--Gwilliam Vol.~I Definition~$6.1.0.5$, adapted to
$\mathrm{Stab}$).
\end{definition}

\begin{remark}[``Locally-defined'' is via curve-stalks]
\label{gluing:rem:quasi-nccr-locally-defined}
The tilting object $T_\sigma$ is defined chamber-by-chamber on
$\mathrm{Stab}(\Db\Coh(X))$; at a generic chamber it is the
K\"unneth tilt $T_{K3, \sigma_1} \boxtimes T_{E, \sigma_2}$ of
the two factor-projected Bridgeland strata, shifted by the Serre
$[3]$. At a wall, the wall-crossing functor $\Psi_W$ updates the
tilt by a spherical twist on the K3 factor or a Poincar\'e-FM twist
on the elliptic factor; the Serre-equivariance axiom ensures the
update commutes with $[3]$.
\end{remark}

\begin{theorem}[Positive-half character of the quasi-NCCR is $-1/\Phi_{10}$]
\label{gluing:thm:quasi-nccr-positive-half-Phi10}
\ClaimStatusProvedElsewhere
The positive-half character of $\mathscr{L}_{K3 \times E}$ equals the
reduced Donaldson--Thomas partition function of $K3 \times E$:
\[
 \chi^+(\mathscr{L}_{K3 \times E})(q, t, p) \;=\;
 Z^{\mathrm{red}}_{\mathrm{DT}}(K3 \times E; q, t, p) \;=\;
 -\frac{1}{\Phi_{10}(q, t, p)} \;=\; -\frac{1}{64\, \Delta_5(q, t, p)^2},
\]
the reciprocal of the Igusa cusp form of weight $10$. The first
equality is Oberdieck--Pandharipande 2016 (arXiv:1603.00276
Theorem~2, proved for primitive curve classes via reduced virtual
classes on $K3$-fibred CY$_3$); the second is the Gritsenko--Nikulin
square-root identity
\[
 \Phi_{10} \;=\; 64\, \Delta_5^2
\]
(Gritsenko--Nikulin 1998 Pt~II arXiv:alg-geom/9810128 Theorem~2.1;
Igusa 1962 \emph{Amer.\ J.\ Math.}~84 on the unique genus-$2$ Siegel
cusp form of weight $10$; Lorgat 2020 arXiv:2004.10817 Theorem~3 for
the additive-lift construction of $\Delta_5$ from $\eta^9 \vartheta_1$).
The Borcherds-lift weight of $\Delta_5$ itself is $\kBKM(\Delta_5) =
c_1(0)/2 = 5$ (Borcherds 1995 \emph{Invent.\ Math.}~120
Theorem~13.3). The numerical factor $64 = 2^6$ records the
$(\Z/2)^6$-character twist of the paramodular additive lift
relative to the Siegel $\Phi_{10}$ normalisation.
\end{theorem}

\begin{proof}
Theorem~\ref{thm:cy-c-beyond-k3e-quasi-nccr-positive-half-character} of
Chapter~\ref{ch:cy-c-beyond-k3e}.
\end{proof}

\subsection{Conditional Hall--Drinfeld double assembly}
\label{gluing:sec:k3xe-hall-drinfeld}

The $24$ curve-stalks, the Mukai lattice cocycle, the Borcherds lift,
and the $M_{24}$ twining supply the ingredients expected to fuse into
one conditional Hopf-algebraic structure. Each ingredient contributes a
well-specified piece: the curve-stalks provide the $24$ elliptic
Hall--positive halves, the Mukai lattice pins down the signature of the
Manin pair, the Borcherds lift supplies the automorphic cocycle, and
the Mathieu twining generates the character lifts. The combined
Hall--Drinfeld double is constructed only after the positive half,
non-degenerate Hopf pairing, completion, bracket comparison, and centre
compatibility have been supplied.

\begin{definition}[Elliptic Hall positive half per curve-stalk]
\label{gluing:def:elliptic-hall-positive-half}
\ClaimStatusDefinition
At each $b \in B^{\mathrm{sing}}$, the local curve-stalk CoHA of
Definition~\ref{gluing:def:local-CoHA-curve-stalk} has a positive-half
summand
\[
 \Yplus_E(b) \;:=\;
 \Yplus(\widehat{\fg\fl}_1) \otimes_\C H^*(E, \C),
\]
an elliptic enhancement of the Schiffmann--Vasserot $\C^3$-CoHA
(Schiffmann--Vasserot 2013 Theorem~1.1), with Hall product twisted
by the Picard--Lefschetz Dehn twist $T_b$.
\end{definition}

\begin{proposition}[Positive-half assembly of $24$ elliptic curve-stalks]
\label{gluing:prop:24-assembly}
\ClaimStatusProvedHere
Fix the local curve-stalk CoHA of
Definition~\ref{gluing:def:local-CoHA-curve-stalk} at every
$b\in B^{\mathrm{sing}}$ and the Picard--Lefschetz transport
$T_b$ of Proposition~\ref{gluing:prop:picard-lefschetz}. For the
chosen DWR/Ran atlas \(\mathfrak U\), this data produces the
curve-stalk positive-half factorisation cosheaf
\[
\Yplus_{\mathrm{stalk}}(K3\times E;\mathfrak U)
 :=
 \mathrm{hocolim}_{\mathrm{Ran}(\bigsqcup_{b\in B^{\mathrm{sing}}}N_b)}
 \{\Yplus_E(b),T_b\}_{b\in B^{\mathrm{sing}}}
\]
with the monodromy relation $\prod_b T_b=1$ imposed on the
$B^{\mathrm{sing}}$ index. This is a positive-half object on the
union of the $24$ elliptic curve-stalks. It is not yet an oriented
compact critical CoHA on all of $K3\times E$.
\end{proposition}

\begin{proof}
Apply the Costello--Gwilliam factorisation-homology machinery
(Vol.~I Definition~$6.1.0.5$, Weiss-cover locality) to the
locally-trivial family $\{\Yplus_E(b)\}_{b\in B^{\mathrm{sing}}}$ on
the disjoint union of curve-stalks. The Ran-space homotopy colimit
computes factorisation sections on disjoint finite configurations in
the $24$ elliptic components, while the parabolic $\SL_2(\Z)$-action
by $T_b$ supplies the descent cocycle in the K3-base direction. The
relation $\prod_b T_b=1$ is precisely the global Picard--Lefschetz
relation on the elliptic K3 base, so the descent datum is consistent.
This constructs the stated factorisation cosheaf. The proof uses no
orientation datum, no compact Hall integration map, no Hopf pairing, and
no Hall--Borcherds bracket comparison; those structures are separate
CY-C inputs.
\end{proof}

\begin{remark}[No Drinfeld double through the homotopy colimit]
\label{gluing:rem:no-double-through-hocolim}
Proposition~\ref{gluing:prop:24-assembly} constructs only the
positive-half object. The expression
\(\mathcal{D}_{\hbar}(\Yplus_{\mathrm{stalk}}(K3\times E;\mathfrak U))\)
is not defined by applying a Drinfeld double to the displayed homotopy
colimit. Such a passage requires a completed monoidal category in
which the positive half is dualizable, a non-degenerate Hopf pairing
compatible with the Picard--Lefschetz cocycle, a bracket comparison
with the Borcherds positive half, and a centre/half-braiding
compatibility. These are the Hall--Drinfeld assembly obligations of
CY-C.
\end{remark}

\begin{proposition}[What the DWR/Ran atlas proves]
\label{gluing:prop:k3xe-ran-atlas-proved-surface}
\ClaimStatusProvedHere
For $X=K3\times E$, the gluing atlas of this section proves exactly the
following Hall-side surface.
\begin{enumerate}[label=\textup{(\roman*)}]
\item It constructs the curve-stalk positive-half cosheaf
\(\Yplus_{\mathrm{stalk}}(K3\times E;\mathfrak U)\) of
Proposition~\ref{gluing:prop:24-assembly}.
\item It identifies the monodromy data that any compact Hall cosheaf
must extend: the Picard--Lefschetz cocycle, the product
$\Aut(K3)\times\Aut(E)$ cocycle, the $M_{24}$ inertia cocycle, and the
Borcherds $\Delta_5$ cocycle.
\item It supplies no morphism
\(\Theta_{\hCS\to\Hall}^{\mathrm{or}}\). Such a morphism exists only
after the oriented hCS--Hall comparison datum and the five vanishing
conditions of Definition~\ref{def:hcs-hall-descent-obstruction} are
supplied, as in Theorem~\ref{thm:hcs-hall-descent-criterion}.
\item It supplies no completed Drinfeld double and no map
\(\Psi^+_{\Hall\to\BKM}\). The double requires the pairing,
completion, centre, and associator data named in
Remark~\ref{gluing:rem:no-double-through-hocolim}; the positive-half
Hall--Borcherds map also requires a Hall product and primitive BPS
bracket compatible with the Borcherds cone of
\(\fg_{\Delta_5}\).
\end{enumerate}
Thus the atlas is a closed curve-stalk positive-half construction and a
complete list of compact comparison obligations; it is not a proof that
\(\CoHA_{\mathrm{crit}}(K3\times E)\), \(D_\hbar\), or
\(G(K3\times E)\) has been constructed.
\end{proposition}

\begin{proof}
Item~\textup{(i)} is Proposition~\ref{gluing:prop:24-assembly}. Item
\textup{(ii)} is the four-stratum decomposition of
Theorem~\ref{gluing:thm:k3xe-four-stratum} together with
Definition~\ref{gluing:def:M24-twining} and
Proposition~\ref{gluing:prop:Delta5-associator}. Item~\textup{(iii)}
is exactly the distinction made by
Definition~\ref{def:cy3-oriented-hcs-hall-comparison-datum} and
Theorem~\ref{thm:hcs-hall-descent-criterion}: a DWR/Ran cover is only
the site on which a comparison datum may live. Item~\textup{(iv)}
follows from the definition of a Drinfeld double of a paired completed
bialgebra and from the Borcherds-side target theorem
Theorem~\ref{thm:k3e-bkm-finite-core}. None of those structures is a
formal consequence of the homotopy colimit defining
\(\Yplus_{\mathrm{stalk}}\).
\end{proof}

\begin{remark}[Finite-height obstruction to double assembly]
\label{gluing:rem:k3xe-finite-height-double-obstruction}
For a height bound \(H\), let
\(\mathfrak O_H(K3\times E)\) denote the ordered obstruction tuple
formed by: the cokernel of the primitive Hall-to-Borcherds multiplicity
map; the radical of the finite-height positive--negative Hall pairing;
the difference between the Hall coproduct and the primitive enveloping
coproduct on associated gradeds; the excess of the Hall kernel over the
Borcherds--Serre ideal; the extra primitive centre; and the mismatch
between the finite-height associator/completion transition maps. Then
\[
 \mathfrak O_H(K3\times E)=0\quad\text{for all }H
\]
is exactly the finite-height input of
Corollary~\ref{cor:k3e-finite-height-promotion-obstruction}. The
DWR/Ran atlas proves the curve-stalk positive half and the monodromy
cocycles; it does not prove the vanishing of these obstruction tuples.
Thus a failure of \(\mathfrak O_H\) at one height blocks both the
completed Hall--Drinfeld double and the BKM current-object promotion,
even though the denominator character \(-\Phi_{10}^{-1}=-\Delta_5^{-2}\)
and \(\kBKM(\Delta_5)=5\) remain theorem-grade.
\end{remark}

\begin{definition}[Mukai-lattice Manin pair on $K3 \times E$]
\label{gluing:def:mukai-manin-pair}
\ClaimStatusDefinition
The \emph{Mukai lattice} of K3 is
\[
 \Lambda_{K3} \;:=\; H^\bullet_{\text{even}}(K3, \Z)
 \;\cong\; \mathrm{II}_{4,20},
\]
a non-degenerate even unimodular lattice of signature $(4, 20)$
(Mukai 1987 \S 3). The Mukai pairing on $\Lambda_{K3}$,
$\langle v, w \rangle = v_0 w_2 + v_2 w_0 - v_1 \cdot w_1$,
provides the (split-signature) Manin-pair datum
$(\sg \oplus \sg^*, \sg, \sg^*)$ needed for a Hall--Drinfeld double
once a positive half, non-degenerate Hopf pairing, completion, and
bracket comparison have been constructed. On $K3 \times E$, the
elliptic factor supplies the chiral central-extension parameter rather
than a new Mukai Cartan summand; the Hall--Drinfeld Manin-pair
signature remains the K3 signature $(4,20)$ in the convention of
Remark~\ref{gluing:rem:signature-conventions}.
\end{definition}

\begin{remark}[The Manin-pair signature is $(4, 20)$ on the K3 factor]
\label{gluing:rem:signature-conventions}
The Manin pair of $\cD_\hbar(K3 \times E)$ is carried by the K3
Mukai lattice alone, signature $(4, 20)$. The elliptic factor
$H^\bullet(E, \Z) = \Z \oplus \Z^2 \oplus \Z$ does not increment
the lattice rank: its rank-$2$ middle cohomology is symplectic (odd
under the elliptic-curve involution and therefore antisymmetric
under the Manin pairing), so it enters the Hall--Drinfeld
construction as the central-extension parameter of the chiral
bracket, not as a new Cartan generator. The resulting imaginary
Cartan count ($23$) is derived in
\S\ref{gluing:sec:k3xe-imaginary-cartan},
Proposition~\ref{gluing:prop:cartan-count}.
\end{remark}

\begin{theorem}[K3 Yangian $Y^{K3} = Y_\hbar(\mathfrak{so}(4, 20))$ and its Hodge-parity $\Z/2$-super-extension]
\label{gluing:thm:k3-super-yangian-osp}
\ClaimStatusConditional{The orthogonal Yangian lane is theorem-grade
given the Hall--Drinfeld comparison input; the Hodge-parity
$\Z/2$-super-extension remains conjectural.}
The Hall--Drinfeld $\hbar$-deformation of the K3 factor is the
orthogonal Yangian
\[
 Y^{K3} \;=\; Y_\hbar(\mathfrak{so}(4, 20))
\]
attached to the Mukai form of signature $(4, 20)$
\textup{(}Molev--Ragoucy~\cite{MolevRagoucy2008} indefinite-orthogonal
Yangian\textup{)}, or conjecturally its Hodge-parity
$\Z/2$-super-extension $Y_\hbar(\mathfrak{so}(4 \mid 20))$ supported by
the Hodge involution $\theta$ on $V = V_+ \oplus V_-$ with
$\dim V_+ = 4$, $\dim V_- = 20$ \textup{(}symmetric on both graded
parts, programme-specific, non-Kac\textup{)}. Kac
$\mathfrak{osp}(4 \mid 20)$ is \emph{not} the envelope: its defining
form is symplectic on the $20$-dimensional odd part, incompatible
with the symmetric Mukai pairing. The $R$-matrix
\[
 R^{K3}(u) \;\in\; \End\bigl(V \otimes V\bigr) \otimes \C(\!(u)\!)
\]
acts on $V = V_+ \oplus V_- = \C^4 \oplus \C^{20}$ with the
Zamolodchikov rational $\mathfrak{so}(N)$ $R$-matrix at $N = 24$,
crossing shift $\kappa_{\so} = N - 2 = 22$, restricted to the
symmetric pair $\mathfrak{so}(4, 20) \supset \mathfrak{so}(4) \oplus
\mathfrak{so}(20)$. The $(4, 20)$-rank is forced by the Mukai form
signature decomposition along the Hodge structure:
the Mukai pairing restricts to signature $(1, 1)$ on $H^0 \oplus H^4$
(hyperbolic plane), $(2, 0)$ on $H^{2, 0} \oplus H^{0, 2}$
(positive-definite, $\R$-rank $2$), and $(1, 19)$ on
$H^{1, 1}_\R$ (K\"ahler class positive, orthogonal $20 - 1 = 19$
negative), summing to $(4, 20)$. The Hodge involution $\theta$ pairs
the four-dimensional $c_+$-part (signature $+$ lines) against the
twenty-dimensional $c_-$-part (signature $-$ lines); on both graded
parts the Mukai pairing remains symmetric, which is the feature that
selects the orthogonal envelope $\mathfrak{so}(4, 20)$ (and its
Hodge-parity $\Z/2$-super-extension $\mathfrak{so}(4 \mid 20)$) and
excludes Kac $\mathfrak{osp}(4 \mid 20)$.
\end{theorem}

\begin{proof}
The Mukai pairing on $\Lambda_{K3} \otimes \R$ is symmetric of
signature $(4, 20)$ and stabiliser $O(4, 20)$; its preserving Lie
algebra is $\mathfrak{so}(4, 20)$ of type $D_{12}$. The
Molev--Ragoucy~\cite{MolevRagoucy2008} indefinite-orthogonal Yangian
$Y_\hbar(\mathfrak{so}(4, 20))$ is the ungraded envelope; the
Hodge-parity $\Z/2$-super-extension
$Y_\hbar(\mathfrak{so}(4 \mid 20))$ is a programme-specific refinement
in which $V_+ \oplus V_-$ carries the Hodge-parity $\Z/2$-grading
and the defining form remains symmetric on both parts (distinct from
Kac $\mathfrak{osp}(4 \mid 20)$, whose odd form is symplectic). The
signature $(4, 20)$ is pinned by the Mukai form signature decomposition
along the K$3$ Hodge structure,
$(1, 1) \oplus (2, 0) \oplus (1, 19)$ on
$(H^0 \oplus H^4) \oplus (H^{2, 0} \oplus H^{0, 2}) \oplus H^{1, 1}_\R$,
summing to $(c_+, c_-) = (4, 20)$. The $4$-dimensional $c_+$-part
decomposes as the K\"ahler line of $H^{1, 1}_\R$, the real-rank-$2$
transcendental plane $H^{2, 0} \oplus H^{0, 2}$, and the
$+$-line of $H^0 \oplus H^4$; the $20$-dimensional $c_-$-part is
the orthogonal complement.
The $R$-matrix $R^{K3}(u)$ is the Maulik--Okounkov stable-envelope
$R$-matrix constructed on the rank-$24$ Mukai space (Maulik--Okounkov
2019 \emph{Ast\'erisque}~408 \S 4), specialised to the symmetric
pair $\mathfrak{so}(4, 20)$; it satisfies the Yang--Baxter equation
on $V^{\otimes 3}$. The conjectural status reflects that the full
K$3$ CoHA-to-orthogonal-Yangian comparison theorem remains open in
the cohomological-Hall literature (cf.\
\S\ref{sec:gluing:obstructions}).
\end{proof}

\begin{remark}[Rigidity at $H\!H^2$: $E_2$-centre obstruction]
\label{gluing:rem:k3-yangian-HH2-rigidity}
The $E_2$-centre $Z_{E_2}(\mathrm{Perf}(K3))$ has $H\!H^2$-rigidity
obstructing any non-abelian deformation of a candidate non-abelian
K$3$ Yangian outside the orthogonal branch
$Y_\hbar(\mathfrak{so}(4, 20))$ and its Hodge-parity
$\Z/2$-super-extension $Y_\hbar(\mathfrak{so}(4 \mid 20))$
(programme-specific, non-Kac): the symmetric structure of the Mukai
pairing on both Hodge-graded parts forces the $\mathfrak{so}$-type,
not the pure-$\mathfrak{gl}$-type nor the Kac orthosymplectic type
$\mathfrak{osp}(4 \mid 20)$ (which would require a symplectic form
on the odd part). See Chapter~\ref{ch:k3-yangian} for the full
$H\!H^2$ analysis.
\end{remark}

\begin{proposition}[$\Delta_5$ Borcherds cocycle as associator datum]
\label{gluing:prop:Delta5-associator}
\ClaimStatusConditional{Requires the compact $K3\times E$ Hall--Drinfeld double: positive half, Hopf pairing, completion, bracket comparison, and centre compatibility.}
The Borcherds-lift cocycle of stratum~(iv) in
Theorem~\ref{gluing:thm:k3xe-four-stratum} supplies the associator
datum expected to endow the completed Hall--Drinfeld double
\(\cD_\hbar(K3 \times E)\), when constructed, with a quasi-Hopf
associator
\[
 \Phi_{\Delta_5} \;\in\;
 \cD_\hbar(K3 \times E)^{\otimes 3}
 \bigl[\![\Delta_5^{-1}]\!\bigr]
\]
with $\kBKM(\Delta_5) = 5$. The associator is the Drinfeld--Kohno
solution of the Knizhnik--Zamolodchikov equation on the Siegel
domain, reduced to the paramodular locus $K(1) \subset \overline{\cA_2}$
where $\Delta_5 \neq 0$. Away from the Humbert loci $H_1, H_4$ the
associator is smooth; crossing a Humbert wall produces an order-$8$
(resp.\ order-$16$) monodromy on $\Phi_{\Delta_5}$.
\end{proposition}

\begin{proof}
Drinfeld 1990 (\emph{Leningrad Math.\ J.}~1, \emph{Quasi-Hopf
algebras}, Theorem~1) constructs quasi-Hopf associators from
$3$-cocycles valued in solutions of the KZ equation. The Siegel-KZ
extension to $\overline{\cA_2}$ is Etingof--Schiffmann--Varchenko
2000 (\emph{Transform.\ Groups} 6 \S 3) for the dynamical twist;
restricting to the paramodular group gives an associator whose
zero locus matches that of $\Delta_5$ on $\overline{\cA_2}$, by
the Bruinier 2002 (Springer LNM 1780) Heegner-Chern-class
reciprocity. The weight $5$ of the associator (i.e.\ the monomial
scaling under the paramodular $\mathbb{G}_m$-action on
$\overline{\cA_2}$) matches $\kBKM(\Delta_5) = c_1(0)/2 = 5$
(Borcherds 1995 Theorem~13.3). This proves the Borcherds cocycle
datum on the paramodular locus. The final step from cocycle datum to
an associator on \(\cD_\hbar(K3\times E)\) is exactly the conditional
double-assembly hypothesis in the statement.
\end{proof}

\begin{definition}[$M_{24}$-twining generating function]
\label{gluing:def:M24-twining}
\ClaimStatusDefinition
For each conjugacy class $[g] \subset M_{24}$, the
\emph{$M_{24}$-twining} of the $K3 \times E$ CoHA is the
generating function
\[
 Z^{[g]}_{K3 \times E}(q, y) \;=\;
 \Tr_{\cH_{K3}^{[g]} \otimes \cH_E}\!\bigl(q^{L_0 - c/24} y^{J_0}
 \zeta_g\bigr),
\]
where $\cH_{K3}^{[g]}$ is the $[g]$-twined sector of the K3 chiral
Hilbert space, $\zeta_g = \mathrm{diag}(\eta^{k_i}_i)$ encodes the
$[g]$-conjugacy class decomposition of $24 = \sum_i k_i$, and $J_0$
is the elliptic $R$-charge. The Eguchi--Ooguri--Tachikawa conjecture
(2010, \emph{Exp.\ Math.}~20; now theorem via Gaberdiel--Hohenegger--
Volpato 2012 and Gannon 2012) identifies each $Z^{[g]}_{K3 \times E}$
with a mock-modular weak Jacobi form $\phi_{[g]} \in J_{0,1}(\Gamma_0([g]))$.
\end{definition}

\begin{remark}[Four cocycles and the conditional Hall--Drinfeld double]
\label{gluing:rem:four-cocycles-one-hopf}
The expected assembly of $\cD_\hbar(K3 \times E)$ is the simultaneous
solution of all four cocycle conditions together with the
positive-half-to-double data:
\begin{itemize}
 \item the Picard--Lefschetz cocycle $\{T_b\}$ of
 Proposition~\ref{gluing:prop:picard-lefschetz} (stratum (i) + (ii));
 \item the $\Aut_s(K3) \times \Aut(E)$-cocycle of the product
 polarisation (stratum (ii));
 \item the $M_{24}$-inertia cocycle of
 Definition~\ref{gluing:def:M24-twining} (stratum (iii));
 \item the $\Delta_5$-associator of Proposition~\ref{gluing:prop:Delta5-associator}
 (stratum (iv)).
\end{itemize}
Compatibility among the four is non-trivial and is the core content of
CY-C. The six routes supply six different construction machines whose
scalar and presentation-level invariants agree on specified loci; their
chiral-algebra convergence is a conjectural bridge/colimit statement,
not a simultaneous six-way isomorphism and not six applications of
\(\Phi_3\). The present section describes the four cocycle ingredients
that the CY-C bridge diagram consumes.
\end{remark}

\subsection{Imaginary rank-23 Cartan}
\label{gluing:sec:k3xe-imaginary-cartan}

The BKM superalgebra $\fg_{\Delta_5}$ whose denominator is $\Delta_5$
has a Cartan subalgebra split into real (finite-root) and imaginary
parts. The imaginary Cartan is the arena where the Heegner
multiplicities live, and its rank is pinned by a short
Mukai-lattice count involving the K3 fibre-line and the elliptic
factor.

\begin{definition}[Imaginary Cartan of $\fg_{\Delta_5}$]
\label{gluing:def:imaginary-cartan}
\ClaimStatusDefinition
The \emph{imaginary Cartan} $\fh^{\mathrm{imag}}_{\Delta_5}$ of
$\fg_{\Delta_5}$ is the sublattice of the weight lattice generated
by the imaginary simple roots $\alpha \in \Lambda_{K3}^{\mathrm{ev}}$
with $\langle \alpha, \alpha\rangle \leq 0$, quotiented by the
Weyl-invariant radical.
\end{definition}

\begin{proposition}[Cartan count]
\label{gluing:prop:cartan-count}
\ClaimStatusProvedHere
The imaginary Cartan of $\fg_{\Delta_5}$ fits in the exact sequence
of $\Q$-vector spaces
\[
 0 \longrightarrow \Q \cdot [\mathrm{fibre}]
 \longrightarrow H^\bullet_{\mathrm{even}}(K3, \Q)
 \longrightarrow \fh^{\mathrm{imag}}_{\Delta_5} \otimes_{\Z} \Q
 \longrightarrow 0,
\]
with the first map sending $1 \mapsto [\text{K3 fibre class}] \in
H^2(K3)$. The rank is
\[
 \rk_\Z\!\bigl(\fh^{\mathrm{imag}}_{\Delta_5}\bigr) \;=\;
 24 - 1 \;=\; 23.
\]
Adjoining the elliptic-curve class $[E] \in H^2(K3 \times E)$ and
again quotienting by the fibre-line redundancy of the elliptic
fibration preserves the count at $23$.
\end{proposition}

\begin{proof}
$\Lambda_{K3} \cong \mathrm{II}_{4,20}$ has rank $24$ over $\Z$
(Mukai 1987 \S 3). The fibre class $[\mathrm{fibre}] \in H^2(K3,
\Z)$ and the zero-section class $[s_0]$ together span a rank-$2$
hyperbolic plane
\[
 U \;=\; \Z[\mathrm{fibre}] \oplus \Z[s_0] \;\subset\; \Lambda_{K3},
 \qquad
 \langle [\mathrm{fibre}], [\mathrm{fibre}]\rangle = 0,\quad
 \langle [s_0], [s_0]\rangle = -2,\quad
 \langle [\mathrm{fibre}], [s_0]\rangle = 1,
\]
(Nikulin 1980 \emph{Izv. Akad. Nauk SSSR Ser. Mat.}~43 Theorem~1.6.1),
so $\Lambda_{K3} = U \oplus U^\perp$ with $U^\perp \cong
\mathrm{II}_{3, 19}$ of rank $22$.

Passing to the imaginary Cartan discards the hyperbolic plane $U$
(real-root redundancy: the two generators of $U$ are the fibre and
section classes, which are real in the Borcherds sense) and retains
$U^\perp \cong \mathrm{II}_{3, 19}$, rank $22$. The Borcherds lift
of the weight-zero weak Jacobi form $\phi_{0, 1} = 2A + 2B$
(Gritsenko--Nikulin 1998 Pt~II, \emph{Int.\ J.\ Math.}~9
Theorem~2.1 constructing $\Delta_5$ as an additive-lift square
root of $\Phi_{10} / 64^2$) adds exactly one imaginary simple root
along the Leech projection of the weight vector (the primitive
isotropic class in $\mathrm{II}_{3, 19}$ transverse to the K3
fibre direction), extending $\mathrm{II}_{3, 19}$ by a single
imaginary generator and lifting the rank to
\[
 \rk_\Z \fh^{\mathrm{imag}}_{\Delta_5} \;=\; 22 + 1 \;=\; 23.
\]
Equivalently, $24 - 1 = 23$: of the $24$ Mukai-lattice generators,
exactly one (the fibre-line redundancy) is absorbed by the Weyl
quotient. The elliptic class $[E] \in H^2(K3 \times E, \Z)$ enters
as the central-extension parameter of the Manin pair
(Definition~\ref{gluing:def:mukai-manin-pair}), not as a new Cartan
generator; its fibre-line redundancy in the elliptic fibration
contributes no rank.
\end{proof}

\begin{remark}[Why not $24$, why not $22$]
\label{gluing:rem:imaginary-cartan-rank}
The Mukai lattice $\Lambda_{K3}$ itself has rank $24$, but this
does not equal the imaginary-Cartan rank: the fibre-line degeneracy
reduces it by $1$ (fibre-line is null in the pairing), giving $23$.
The further quotient by the hyperbolic plane $U$ would give $22$,
but that quotient is too aggressive; one loses the imaginary
simple root along the Leech direction that the Borcherds lift of
$\phi_{0,1}$ adds back. The count $23$ is the self-consistent
Heegner-multiplicity count with $\Delta_5$'s Fourier coefficient at
$(N, \ell) = (0, 0)$ supplying $c_1(0) = 10$ and the weight-$5$
specialisation.
\end{remark}

\begin{corollary}[Imaginary simple roots seeded by $\Delta_5$-Fourier coefficients]
\label{gluing:cor:imaginary-seeded}
\ClaimStatusProvedElsewhere
Each imaginary simple root $\alpha \in \fh^{\mathrm{imag}}_{\Delta_5}$
with $\langle \alpha, \alpha\rangle = 4N - \ell^2 \leq 0$ has
multiplicity $c_{\Delta_5}(N, \ell)$, the Fourier coefficient of
$\Delta_5$ on $\overline{\cA_2}$ at the Jacobi parameter $(N, \ell)$.
\end{corollary}

\begin{proof}
This is Theorem~\ref{thm:k3e-denominator} of
Chapter~\ref{ch:k3e-bkm} (the Gritsenko--Nikulin denominator identity
with imaginary-root Fourier coefficients). The weight-$5$ of
$\Delta_5$ is the Borcherds BKM weight, $\kBKM = 5$.
\end{proof}

\subsection{Humbert monodromy and $\kch + \kch^! = 8$}
\label{gluing:sec:k3xe-humbert}

The Kuga--Satake period map carries $K3 \times E$ into the
paramodular Siegel moduli $\overline{\cA_2}$. The Borcherds form
$\Delta_5$ vanishes on two Humbert divisors $H_1, H_4 \subset
\overline{\cA_2}$; the monodromy of the holonomic
$\cD$-module $\cL^{\Delta_5}$ around each divisor has a specific
order. The Humbert-$H_1$ order equals $8$; this is the $\B$-row
Koszul conductor in the derived-centre complementarity
$\kch + \kch^! = 8$ of Vol~I Theorem~C.

\begin{proposition}[Kuga--Satake abelian variety of K3: Deligne dimension]
\label{gluing:prop:kuga-satake-dimension}
\ClaimStatusProvedHere
Let $T(K3)_{\Q} \subset H^2(K3, \Q)$ be the transcendental lattice
of a (generic) K3 surface, of rank $n = 22$. The Kuga--Satake
abelian variety $\mathrm{KS}(K3)$ associated to the weight-two
Hodge structure on $T(K3)_\Q$ has complex dimension
\[
 \dim_\C \mathrm{KS}(K3) \;=\; 2^{n - 2} \;=\; 2^{20}.
\]
For a lattice-polarised K3 with Picard rank $\rho$, the
transcendental rank is $n = 22 - \rho$ and
$\dim_\C \mathrm{KS} = 2^{20 - \rho}$.
\end{proposition}

\begin{proof}
Deligne 1972 (\emph{La conjecture de Weil pour les surfaces K3},
\emph{Invent.\ Math.}~15, Proposition~4.5) constructs the
Kuga--Satake abelian variety from a polarised weight-two
$\Q$-Hodge structure $(V, \langle -, - \rangle)$ of signature
$(2, n - 2)$ by extracting the even Clifford algebra
$C^+(V, \langle -, - \rangle)$, a semisimple $\Q$-algebra of
dimension $2^{n - 1}$. The Hodge structure on $V$ induces a
complex structure on $C^+(V) \otimes_\Q \R$ whose kernel lattice
is a polarisable abelian variety $\mathrm{KS}$; the complex
dimension of $\mathrm{KS}$ is half the real rank of $C^+(V) \otimes \R$,
giving $\tfrac12 \cdot 2^{n - 1} = 2^{n - 2}$ (Deligne 1972
Proposition~4.5; the calculation is the Clifford-algebra half-dimension
count, independent of the signature of $\langle -, - \rangle$).
For a K3 surface the total rank is $\rk H^2 = 22$; on the generic
K3 the transcendental lattice $T(K3)$ saturates this, giving
$n = 22$, $\dim_\C \mathrm{KS} = 2^{20}$. For lattice-polarised
K3 with Picard rank $\rho$, the transcendental rank is $22 - \rho$,
yielding $2^{20 - \rho}$.
\end{proof}

\begin{remark}[Kuga--Satake period map to $\overline{\cA_2}$]
\label{gluing:rem:ks-map-A2}
The Kuga--Satake abelian variety $\mathrm{KS}(K3)$ is a $2^{20}$-
dimensional abelian variety whose moduli space receives the K3
period map. For \emph{polarised} K3 with transcendental rank $n = 3$
(the narrow rank relevant to $\overline{\cA_2}$: two-dimensional
principally polarised abelian varieties), $\dim_\C \mathrm{KS} =
2^{3 - 2} = 2$, identifying $\mathrm{KS}$ with a principally polarised
abelian surface and placing the period point in the paramodular
moduli $\overline{\cA_2} = \cA_2 / \mathrm{para}$ on which
$\Delta_5$ is a Borcherds form (Gritsenko--Nikulin 1998 Pt~II
Theorem~2.1). The $K3 \times E$ threefold realises this restriction:
the product structure imposes a lattice-polarisation of rank
$\rho = 19$ on the K3 factor (K\"ahler $+$ elliptic fibre $+$ section
$+$ 16 rational nodal components), leaving transcendental rank
$22 - 19 = 3$ and $\mathrm{KS}$-dimension $2$; a principally
polarised abelian surface as required.
\end{remark}

\begin{definition}[Humbert surface $H_1 \subset \overline{\cA_2}$]
\label{gluing:def:humbert-H1}
\ClaimStatusDefinition
The \emph{Humbert surface} $H_1 \subset \overline{\cA_2}$ is the
Heegner divisor of discriminant $1$: the locus of principally
polarised abelian surfaces that split as a product of two elliptic
curves $E_1 \times E_2$ with the product polarisation. Equivalently,
$H_1$ is the zero locus of the simplest non-trivial Humbert
modular function of discriminant $1$ on $\overline{\cA_2}$. The
paramodular form $\Delta_5$ vanishes on $H_1$ to order $1$ on
generic coordinates, order $2$ on $H_4$.
\end{definition}

\begin{theorem}[Humbert-$H_1$ monodromy order is $8$]
\label{gluing:thm:humbert-order-8}
\ClaimStatusProvedElsewhere
The local monodromy of the holonomic $\cD$-module $\cL^{\Delta_5}$
around $H_1 \subset \overline{\cA_2}$ has order $8$. Equivalently, on
the conditional Hall--Drinfeld assembly locus, the Koszul conductor of
the Mukai-doubled K3 positive-half/Manin-pair datum satisfies
\[
 K^{\kch}(K3 \times E) \;=\; 2\, c_+(\mathrm{Mukai}(K3))
 \;=\; 2 \cdot 4 \;=\; 8.
\]
\end{theorem}

\begin{proof}
The Mukai pairing on $H^\bullet_{\mathrm{even}}(K3, \R)$ decomposes
along the K3 Hodge structure as
\[
 \bigl(H^0 \oplus H^4\bigr) \;\oplus\; \bigl(H^{2, 0} \oplus H^{0, 2}\bigr)
 \;\oplus\; H^{1, 1}_{\R},
\]
with Mukai-signature components $(1, 1) \oplus (2, 0) \oplus (1, 19)$
(the middle factor $H^{2, 0} \oplus H^{0, 2}$ is positive-definite of
real rank $2$; the K3 Hodge index pins the $(1, 19)$ split of
$H^{1, 1}_{\R}$; the extremal pair $H^0 \oplus H^4$ is hyperbolic).
The total $(4, 20)$ is the sum; $c_+ = 1 + 2 + 1 = 4$ (Mukai 1987
\S 3; Huybrechts 2016 \emph{Lectures on K3 Surfaces} Ch.~14 \S 2;
the decomposition is the Manin correction recorded in the
adjudication ledger for the Mukai-Hodge split). Bruinier's
Heegner-divisor Chern-class reciprocity (Bruinier 2002 Springer
LNM 1780, Proposition~5.1 and the Borcherds product theorem \S 3.2)
identifies the local monodromy of a Borcherds lift around a Heegner
divisor $Z(m, \mu)$ as the image in $\SL_2(\Z / 2N\Z)$ of the
Siegel--Weil reciprocity generator, with order
\[
 \mathrm{ord}_{Z(m, \mu)}(\cL^\Psi) \;=\; 2\, c_+\, \nu(\mu)
\]
where $\nu(\mu)$ is the $\mu$-torsion order in the discriminant
group of $\Lambda_{K3}$ (Bruinier 2002 Theorem~5.6).

For $\Psi = \Delta_5$ and the Humbert divisor $H_1 = Z(1, \mu_0)$
with $\mu_0$ the primitive generator of the discriminant form
(discriminant $1$, $\nu(\mu_0) = 1$), one reads
\[
 \mathrm{ord}_{H_1}(\cL^{\Delta_5}) \;=\; 2 \cdot c_+ \cdot 1
 \;=\; 2 \cdot 4 \;=\; 8.
\]
The Gritsenko additive lift $\eta^9 \vartheta_1$ of weight $9/2$ and
index $1/2$, with Fourier coefficient $c_{\eta^9 \vartheta_1}(1) =
\pm 1/4$ at norm $1$ (Gritsenko--Nikulin 1998 Pt~I,
\emph{Int.\ J.\ Math.}~9 \S 2 Table~1; Lorgat 2020
arXiv:2004.10817 Theorem~3), provides the explicit input:
the $\Z/8$-multiplier system on the paramodular $\mathbb{Z}/2$-spin
cover is generated by $(\eta^9 \vartheta_1)^8 \cdot \cO(-H_1)^{-1}$
trivialising to $\cO_{\overline{\cA_2}}$. The Riemann--Hilbert
correspondence (Deligne 1970 \emph{\'Equations diff\'erentielles
\`a points singuliers r\'eguliers} LNM 163 Theorem~II.1.9;
Kashiwara 1984 \emph{Publ.\ RIMS}~20) translates torsion order of
the associated $\cD$-module to monodromy order, giving $8$.
\end{proof}

\begin{remark}[The $8$ as $2 c_+$ of the Mukai signature]
\label{gluing:rem:eight-from-signature}
The integer $8$ is not a coincidence. The Mukai lattice
$\mathrm{II}_{4, 20}$ carries signature $(1, 1) \oplus (2, 0) \oplus
(1, 19)$ on $H^0 \oplus H^4$, $H^{2,0} \oplus H^{0, 2}$, and
$H^{1, 1}_{\R}$ respectively, summing to $(c_+, c_-) = (4, 20)$ with
$c_+ = 1 + 2 + 1 = 4$; the Mukai-doubling that produces the
Hall--Drinfeld pair $(\sg, \sg^*)$ from the self-dual
$\Lambda_{K3}$ multiplies by $2$, and $2 \cdot c_+ = 2 \cdot 4 = 8$
is the Koszul conductor $K^{\kch}$. The other summand $c_- = 20$
enters separately in the split $\kfib(K3) = c_+ + c_- = 24$,
corresponding to the separation between $\kch + \kch^! = 8$ (the
$\B$-row complementarity of Vol~I Theorem~C) and the Mukai-rank
$\kfib(K3) = 24$.
\end{remark}

\begin{theorem}[Derived-centre complementarity on the $\B$-row]
\label{gluing:thm:B-row-complementarity}
\ClaimStatusProvedElsewhere
The $K3 \times E$ landmark of the canonical five-archetype
$(\mathsf G, \mathsf L, \mathsf C, \mathsf M, \mathsf B)$ landmark
ceiling (Vol~I Theorem~C) satisfies
the complementarity identity
\[
 \kch + \kch^{!} \;=\; 8 \quad\text{(on the}\ \mathsf B\text{-row)},
\]
where $\kch^{!}$ is the Verdier-dual Koszul datum and $8 = K^{\kch}$
is the Humbert-$H_1$ monodromy order of
Theorem~\ref{gluing:thm:humbert-order-8}.
\end{theorem}

\begin{proof}
Bruinier Heegner-Chern-class reciprocity identifies $K^{\kch} = 8$ as
both the monodromy order and the Koszul conductor; the
complementarity $\kch + \kch^! = K^{\kch}$ is Vol~I Theorem~C on the
five-archetype ceiling, specialised to the $\mathsf B$-row via the
Mukai-enhanced K3 Heisenberg witness. See
Remark~\ref{rem:mukai-doubling-K-kappa} of
Chapter~\ref{ch:k3-chiral-bialgebra-platonic} for the full argument.
\end{proof}

\begin{corollary}[$H_4$ doubles $H_1$: order $16$]
\label{gluing:cor:H4-order-16}
\ClaimStatusProvedElsewhere
The local monodromy of $\cL^{\Delta_5}$ around $H_4 \subset
\overline{\cA_2}$ has order $16$, twice the Humbert-$H_1$ order,
because $c_{\eta^9 \vartheta_1}(4) = \pm 1/8$ at norm $4$ halves the
denominator and the index-$1/2$ contribution doubles the effective
order.
\end{corollary}

\begin{proof}
As in Theorem~\ref{gluing:thm:humbert-order-8}, with the Fourier
coefficient at norm $4$ (discriminant $4$) producing $1/8$ in place
of $1/4$; doubling via the index-$1/2$ contribution gives $16$.
\end{proof}

\begin{remark}[Stratum (iv) is invisible to Cech gluing]
\label{gluing:rem:iv-invisible-cech}
The Humbert monodromy identity $\kch + \kch^! = 8$ is a
\emph{stratum-(iv)} invariant: it refers to the Borcherds lift on
the Kuga--Satake period domain, not to any Cech atlas on $X$ itself.
No chart-wise (stratum-(i)) or derived-Morita (stratum-(i)$\cap$(ii))
gluing could produce it: the Borcherds form $\Delta_5$ is a global
automorphic section on $\overline{\cA_2}$, not a locally assembled
datum on $X$. This is the reason the master example requires the
full four-stratum atlas: three strata alone cannot produce the
$K^{\kch} = 8$ relation that anchors the derived-centre
complementarity.
\end{remark}

\subsection{Synthesis: why $K3 \times E$ is the master}
\label{gluing:sec:k3xe-synthesis}

The $K3 \times E$ threefold is the \emph{only} compact CY$_3$ in
Vol~III on which all four equivariance strata act simultaneously and
non-trivially. The conifold and local $\PP^2$ live in strata (i)
and (i)$\cap$(iii) respectively; K3 alone lives in (ii)$\cap$(iii);
the quintic lives in none of the four with explicit-formula content
and accesses only the stratum-free categorical CoHA. The four-stratum
simultaneity gives $K3 \times E$:
\begin{itemize}
 \item \emph{Curve-stalk enumeration} (\S\ref{gluing:sec:k3xe-elliptic-fibration},
 \S\ref{gluing:sec:k3xe-local-global-assembly}): exactly $24$
 $I_1 \times E$ curve-stalks via Kodaira $\chi = 24$ on an Ehresmann
 fibre bundle; local model $\{xy = w\} \times E$ at each; Picard--Lefschetz
 parabolic $T_b = \begin{psmallmatrix} 1 & 1 \\ 0 & 1 \end{psmallmatrix}
 \in \SL_2(\Z)$ in the symplectic basis $(\delta, \delta^\vee)$.
 \item \emph{Four-stratum decomposition}
 (\S\ref{gluing:sec:k3xe-four-stratum}): latent toric germ, reduced
 $\Aut$-cocycle, $M_{24}$-inertia, $\Delta_5$-Borcherds.
 \item \emph{Quasi-NCCR substitute}
 (\S\ref{gluing:sec:k3xe-quasi-nccr}): five-fold global-NCCR
 obstruction resolved by the Serre-equivariant tilted-algebra sheaf
 on $\mathrm{Stab}(\Db\Coh(X))$, with $\chi^+(\mathscr{L}_X) =
 -1/\Phi_{10} = -1/(64\,\Delta_5^2)$ the reduced DT partition.
 \item \emph{Hall--Drinfeld assembly}
 (\S\ref{gluing:sec:k3xe-hall-drinfeld}): $24$ elliptic Hall
 positive halves, Mukai Manin pair at signature $(4, 20)$, $\Delta_5$-
 associator for the quasi-Hopf structure, $M_{24}$-twining, and the
 K$3$ Yangian $Y^{K3} = Y_\hbar(\mathfrak{so}(4, 20))$
 (Molev--Ragoucy orthogonal Yangian) with conjectural Hodge-parity
 $\Z/2$-super-extension $Y_\hbar(\mathfrak{so}(4 \mid 20))$ on
 $V = V_+ \oplus V_- = \C^4 \oplus \C^{20}$ (non-Kac; Kac
 $\mathfrak{osp}(4 \mid 20)$ is excluded as its odd form is
 symplectic; Theorem~\ref{gluing:thm:k3-super-yangian-osp}).
 \item \emph{Imaginary-Cartan rank}
 (\S\ref{gluing:sec:k3xe-imaginary-cartan}): rank $23 = 22 + 1 =
 \rk(\mathrm{II}_{3, 19}) + \text{Leech addback}$ with
 imaginary-simple-root multiplicities $c_{\Delta_5}(N, \ell)$.
 \item \emph{Kuga--Satake dimension} and \emph{derived-centre
 complementarity} (\S\ref{gluing:sec:k3xe-humbert}):
 $\dim_\C \mathrm{KS}(K3) = 2^{n - 2}$ (Deligne 1972
 Proposition~4.5; Proposition~\ref{gluing:prop:kuga-satake-dimension});
 Humbert-$H_1$ monodromy order $8 = 2 c_+ \nu(\mu_0)$ with $c_+ = 4$
 from the Mukai Hodge signature $(1, 1) \oplus (2, 0) \oplus (1, 19)$,
 giving $\kch + \kch^! = 8$.
\end{itemize}

\begin{theorem}[Four-$\kappa_\bullet$ spectrum on $K3 \times E$ from four distinct constructions]
\label{gluing:thm:four-kappa-spectrum}
\ClaimStatusProvedElsewhere
On $X = K3 \times E$, the four subscripted $\kappa$-invariants take
the values $\{0, 3, 5, 24\}$. Each value is the output of a
\emph{different} construction; no two values are identifications of
the same object through a single functor, and the tuple is not
$\Phi_3$ applied four times:
\begin{center}
\begin{tabular}{|l|l|l|}
\hline
Construction & Invariant & Value \\ \hline
K\"unneth total space, $\chi(\cO_X) = \chi(\cO_{K3}) \cdot
\chi(\cO_E)$ & $\kcat(K3 \times E)$ & $2 \cdot 0 = 0$ \\
Chiral Heisenberg--Mukai specialisation of $\Phi_3$ at $X$ &
$\kch^{\mathrm{Heis}}(K3 \times E)$ & $\dim_\C X = 3$ \\
Borcherds lift along the Kuga--Satake period map &
$\kBKM(\fg_{\Delta_5})$ & $c_1(0)/2 = 5$ \\
Mukai-lattice rank of the K3 fibre, $\chi_{\mathrm{top}}(K3)$ &
$\kfib(K3)$ & $24$ \\ \hline
\end{tabular}
\end{center}
The four rows inhabit four distinct categorical slots:
$\kcat$ lives in derived global sections of $\cO_X$, $\kch$ in
chiral-algebra factorisation homology, $\kBKM$ in automorphic
Borcherds-product weights, $\kfib$ in Mukai-vector even cohomology.
\end{theorem}

\begin{proof}
Each row is its own cited theorem:
$\kcat(K3 \times E) = 0$ by K\"unneth multiplicativity
$\chi(\cO_{X \times Y}) = \chi(\cO_X) \chi(\cO_Y)$
(Hartshorne 1977 \emph{Algebraic Geometry} Ch.~III Exercise~8.3;
$\chi(\cO_{K3}) = 2$, $\chi(\cO_E) = 0$; recorded in
Chapter~\ref{ch:cy-d-kappa-stratification}
Theorem~\texttt{thm:chi-O-kunneth-total-space});
$\kch^{\mathrm{Heis}}(K3 \times E) = 3$ by the Heisenberg--Mukai
chiral specialisation
(Chapter~\ref{ch:k3-chiral-bialgebra-platonic}
Theorem~\ref{thm:k3-chiral-bialgebra-four-kappas});
$\kBKM(\Delta_5) = 5$ by Gritsenko--Nikulin 1998 Pt~II Theorem~2.1
($c_1(0) = 10$, paramodular weight $5$);
$\kfib(K3) = 24$ by the Noether formula $\chi_{\mathrm{top}}(K3) =
c_2(K3) = 24$ (Proposition~\ref{gluing:prop:kodaira-24}).
The four slots never meet: the K\"unneth identity on $\kcat$ is
multiplicative on products and kills a factor with $\chi(\cO_E) =
0$; the chiral specialisation $\kch^{\mathrm{Heis}} = \dim_\C X$
records the $E_1$-shadow shift-law value at $d = 3$; the Borcherds
weight is an automorphic-theoretic invariant of the Kuga--Satake
image of $X$; the Mukai rank is a lattice invariant of the K3
fibre alone. No two depend on each other functorially.
\end{proof}

\begin{remark}[The attempted reduction $\kBKM = \kch + \chi(\cO_{\mathrm{fibre}})$ fails]
\label{gluing:rem:failed-reduction}
One might hope for a unifying formula $\kBKM = \kch +
\chi(\cO_{\mathrm{fibre}})$; at $N = 1$, the right-hand side is
$\kch(K3 \times E) + \chi(\cO_E) = 0 + 0 = 0$, not $5$. At $N = 2$,
the right-hand side is $1$, the left-hand side is $4$. The universal
formula is $\kBKM(\Phi_N) = c_N(0)/2$ (Borcherds 1995 Theorem~13.3;
Gritsenko 1999 \emph{St.\ Petersburg Math.\ J.}; Chapter~\ref{ch:cy-d-kappa-stratification}
Theorem~\texttt{thm:borcherds-weight-kappa-BKM-universal}), not a
sum of $\kch$ and a fibre contribution. The four constructions
simply do not reduce to each other.
\end{remark}

\begin{remark}[Master status confirmed]
\label{gluing:rem:master-status}
$K3 \times E$ is the master example of the gluing atlas: every
Part of this chapter's mathematical content is simultaneously
active on this one compact CY$_3$, with each part contributing
exactly one non-redundant piece. The classification theorem of
\S 10.1 lands on $K3 \times E$ in the cell
$(\text{i}) \cap (\text{ii}) \cap (\text{iii}) \cap (\text{iv})$,
with the full nested fibre-product cocycle data from
Definition~\ref{gluing:def:mukai-manin-pair},
Proposition~\ref{gluing:prop:Delta5-associator},
Definition~\ref{gluing:def:M24-twining}, and the $24$ Picard--Lefschetz
curve-stalks of Proposition~\ref{gluing:prop:24-assembly}. No other
compact CY$_3$ in the Vol~III landscape activates all four.
\end{remark}

\paragraph{What was just established, what is next.}
The curve-stalk enumeration placed the local-model obstruction at
the right categorical level (Jordan-triple with elliptic
coefficients, not smooth $\C^3$-point). The four-stratum
decomposition established that $K3 \times E$ sits in the deepest
cell of the equivariance atlas with four logically independent
gluing cocycles. The Serre-equivariant quasi-NCCR substituted for
the absent global NCCR via a tilted-algebra sheaf over
$\mathrm{Stab}(\Db\Coh(X))$. The Hall--Drinfeld assembly obligations,
the imaginary-Cartan count, and the Humbert monodromy identity
pin the conditional route by which $\fg_{\Delta_5}$ would arise from
the positive-half CoHA data, with denominator $\Delta_5$ and
complementarity $\kch + \kch^! = 8$ after the required pairing,
completion, bracket, and centre data are constructed.
\S\ref{sec:gluing:obstructions} addresses the regime where gluing fails: compact
CY$_3$'s without the $K3 \times E$ fibration structure (quintic,
bicubic) for which the four-stratum atlas is unavailable and
only the stratum-free categorical CoHA via MNOP + Davison
integrality applies.
